\section{Monitores e \texttt{synchronized}}
\label{sec:monitores}

\paragraph{O problema.}
Na JLS 2ª ed., as threads executam código que opera sobre valores e
objetos de uma memória principal compartilhada. A linguagem apoia a escrita de
programas que, mesmo concorrentes, tenham comportamento determinístico,
e para isso fornece mecanismos de sincronização (cap.~17, abertura). O
mesmo texto descreve a alternativa a evitar: \enquote{spinning}, travar
e destravar um objeto repetidamente para ver se o estado interno mudou,
o que consome esforço computacional.

\paragraph{A decisão.}
A resposta da JLS 2ª ed. é o \termo{monitor}{Mecanismo que deixa uma
única thread por vez executar a região de código que ele protege.}:
\enquote{a high-level mechanism for allowing only one thread at a time
to execute a region of code protected by the monitor}. O comportamento
do monitor é explicado em termos de \termo{locks}{Travas. Uma thread
adquire o lock antes de entrar na região protegida e o libera ao
sair.}, um por objeto (cap.~17, abertura). A linguagem não oferece
ações separadas de travar e destravar: o bloco \texttt{synchronized} e
o método \texttt{synchronized} as executam aos pares. Um método de
instância trava o lock de \texttt{this}, e um método \texttt{static}
trava o do objeto \texttt{Class} da classe (§17.13). Os métodos
\texttt{wait}, \texttt{notify} e \texttt{notifyAll} de \texttt{Object}
dão suporte à transferência eficiente de controle entre threads: em vez
de girar, uma thread pode se suspender com \texttt{wait} até outra
acordá-la com \texttt{notify}
(cap.~17, abertura).

Cada objeto tem, além do lock, um \termo{wait set}{Conjunto das
threads que chamaram \texttt{wait} no objeto e ainda não foram
removidas dele.} (§17.14). Segundo a §17.14, \texttt{wait} só deve ser
chamado por quem já detém o lock do objeto. A chamada põe a thread no
wait set e executa tantas liberações quantos forem os travamentos
pendentes. A thread fica dormente até acontecer uma de três coisas: um
\texttt{notify} que a escolha arbitrariamente, um \texttt{notifyAll} ou
o fim do tempo limite, se houver. Depois ela trava o objeto de novo e
volta de \texttt{wait} com o lock no mesmo estado em que estava.

A JVMS 2ª ed. traz o modelo de threads e locks no cap.~8, \enquote{adapted with
minimal changes from Chapter 17 of the first edition of The Java
Language Specification}, nas seções §8.12 a §8.14. Na §8.14, um
\texttt{notify} remove do wait set uma thread escolhida arbitrariamente.

Sobre a origem, o white paper diz que as primitivas de
sincronização seguem o paradigma de monitor e variável de condição
\enquote{introduced twenty years ago by C.A.R. Hoare}, e que foram
modeladas no Cedar/Mesa (§7.1 do white paper,
\citetitle{paper:java-language-environment}). O artigo de Hoare
(\citetitle{paper:hoare1974-monitors}, 1974) \enquote{develops
Brinch-Hansen's concept of a monitor as a method of structuring an
operating system}, com variáveis de condição e as operações
\texttt{wait} e \texttt{signal}. Em 1999, Per Brinch Hansen
escreveu em \citetitle{paper:brinch-hansen1999-insecure}:
\enquote{The claim that Java supports monitors is shown to be false.}
Para ele, o que o Java oferece \enquote{is just a programming style that
imitates insecure monitors}, e o artigo conclui que o Java ignora os
últimos vinte e cinco anos de pesquisa em linguagens de programação
paralela.

\paragraph{O que mudou depois.}
Na JLS do Java SE~25, o cap.~17 continua se chamando
\enquote{Threads and Locks}, mas a sincronização está na §17.1, o wait
set na §17.2 e o modelo de memória na §17.4. A §17.1 define:
\enquote{Each object in Java is associated with a monitor, which a
thread can lock or unlock.} Só uma thread por vez detém o lock de um
monitor. A §17.1 descreve o bloco e o método \texttt{synchronized} com o texto
da §17.13 da 2ª edição, exceto que o que se trava é o monitor do
objeto e sai a observação sobre a memória principal. A mesma seção cita as classes de
\texttt{java.util.concurrent} entre os outros meios de sincronização.

A §17.2.1 detalha o \texttt{wait}. Se a thread não detém o lock do
objeto, a chamada lança \texttt{IllegalMonitorStateException}. Se
detém, ela entra no wait set e executa uma liberação para cada
travamento pendente; ao sair, refaz os mesmos travamentos. A lista de saídas do wait set cresceu: além
de \texttt{notify}, \texttt{notifyAll} e tempo limite, entram a
interrupção da thread e uma ação interna da implementação. Essa última é
o \termo{spurious wakeup}{Acordar espúrio: a thread sai de
\texttt{wait} sem \texttt{notify}, interrupção nem fim do tempo
limite.}: \enquote{Implementations are permitted, although not
encouraged, to perform \enquote{spurious wake-ups}}. A mesma seção diz que
essa permissão exige usar \texttt{wait} só dentro de laços que só
terminam quando a condição esperada vale. A §17.2.2 não dá garantia sobre qual
thread o \texttt{notify} escolhe, e, após um \texttt{notifyAll}, só uma
thread por vez retoma o lock. O javadoc de \texttt{Object.wait} no
JDK~25 (\citetitle{doc:javase25-object}) recomenda testar a condição
num laço \texttt{while} em volta da chamada.

Na JVMS do Java SE~25, a §2.11.10 diz que a JVM sincroniza métodos e
sequências de instruções com \enquote{a single synchronization
construct: the monitor}. A sincronização é feita pela entrada e saída
do monitor, explícita, com \texttt{monitorenter} e \texttt{monitorexit},
ou implícita, nas instruções de chamada e retorno de método (§3.14).

\begin{center}
\begin{tabular}{@{}p{0.22\textwidth}p{0.36\textwidth}p{0.36\textwidth}@{}}
\toprule
Aspecto & JLS/JVMS 2ª ed. (§) & SE 25 (§) \\
\midrule
Onde está & JLS §17.12--§17.14; JVMS §8.12--§8.14 & JLS §17.1 e
§17.2; JVMS §2.11.10 e §3.14 \\
Unidade & Um lock por objeto (JLS §17.13) & Um monitor por objeto
(JLS §17.1) \\
\texttt{wait} sem o lock & Só deve ser chamado com o lock (JLS §17.14)
& Lança \texttt{IllegalMonitor\-StateException} (JLS §17.2.1) \\
Saídas do wait set & \texttt{notify}, \texttt{notifyAll}, tempo limite
(JLS §17.14) & Também interrupção e ação interna (JLS §17.2.1) \\
Spurious wakeup & Não mencionado (JLS cap.~17) & Permitido, não
encorajado (JLS §17.2.1) \\
\bottomrule
\end{tabular}
\end{center}

\paragraph{Onde fica o \texttt{wait} dentro do monitor?}
O exemplo abaixo é um buffer de duas posições. \texttt{colocar} espera
enquanto o buffer está cheio, \texttt{retirar} espera enquanto está
vazio, e os dois chamam \texttt{notifyAll} depois de mudar o estado.

\lstinputlisting[language=Java]{generated/source/EsperaEmLaco.java}

\inputbytecode{EsperaEmLaco}

\begin{itemize}
\item O cabeçalho mostra \texttt{synchronized void colocar(int)}, mas o
  código não tem \texttt{monitorenter} nem \texttt{monitorexit}. O
  método leva a flag \texttt{ACC\_SYNCHRONIZED}, conferida pelas
  instruções de chamada: a thread entra no monitor de \texttt{this} antes
  do offset~0 e sai dele no retorno, normal ou por exceção (JVMS SE~25,
  §2.11.10). O bloco \texttt{synchronized}, com as duas instruções, está
  na seção~\ref{sec:sincronizacao}.
\item Em \texttt{colocar}, os offsets 0--9 testam
  \texttt{quantidade == itens.length}; o \texttt{if\_icmpne} do
  offset~9 sai do laço para o offset~19. Nos offsets 12--13,
  \texttt{aload\_0} e \texttt{invokevirtual} chamam
  \texttt{Object.wait} sobre \texttt{this}, o mesmo objeto cujo monitor
  o método detém. O \texttt{goto} do offset~16 volta ao offset~0, e a
  condição é testada de novo.
\item Dentro da chamada do offset~13, a thread executa as liberações
  e, ao sair do wait set, os travamentos de volta, antes de
  \texttt{wait} retornar (JLS SE~25, §17.2.1). Entre as liberações e
  os travamentos de volta, a thread não detém o monitor de
  \texttt{this}: o javadoc de \texttt{Object.wait}
  (\citetitle{doc:javase25-object}) diz que ela passa a
  \enquote{relinquish any and all synchronization claims on this
  object}.
\item Em \texttt{retirar}, o padrão se repete: \texttt{ifne} no
  offset~4, \texttt{wait} no offset~8 e \texttt{goto} para o offset~0
  no offset~11. \texttt{notifyAll} está no offset~51 dos dois métodos,
  antes do \texttt{return} (54) e do \texttt{ireturn} (55).
\end{itemize}

Experimento (JDK~25, imagem \texttt{eclipse-temurin:25-jdk}): a saída
real de \texttt{main} foi \texttt{1}, \texttt{2}, \texttt{3},
\texttt{4} e \texttt{5}, um número por linha.

\paragraph{O que isso muda no seu código.}
Cada recomendação tem seu escopo.
\begin{itemize}
\item \emph{JDK~25:} chame \texttt{wait} dentro de um \texttt{while}
  que testa a condição, nunca de um \texttt{if}. A JLS SE~25 permite o
  spurious wakeup (§17.2.1), e o javadoc de \texttt{Object.wait} do
  JDK~25 recomenda o laço. O javadoc já descreve o spurious wakeup e
  pede o laço desde o Java~5, de 2004 (\citetitle{doc:javase5-object}).
\item \emph{JLS 2ª ed. e JDK~25:} todo objeto tem um lock (JLS 2ª
  ed., §17.13) ou monitor (JLS SE~25, §17.1), que qualquer thread com
  uma referência a ele pode travar. O exemplo acima trava \texttt{this}
  de propósito, para mostrar um método \texttt{synchronized} sem
  \texttt{monitorenter}; a flag \texttt{ACC\_SYNCHRONIZED} só aparece
  com \texttt{javap -v}. Por que preferir um objeto de lock privado está
  na seção~\ref{sec:sincronizacao}.
\item \emph{Java~5 (2004) em diante:} para código novo, prefira as
  classes de \texttt{java.util.concurrent}, que a própria §17.1 da JLS
  SE~25 cita como alternativa de sincronização. O pacote veio da
  \citetitle{doc:jsr166}, com versão final em 30 de setembro de 2004.
\end{itemize}

\paragraph{Onde aprofundar.}
O \texttt{monitorenter} e o \texttt{monitorexit} estão na
seção~\ref{sec:sincronizacao}. As garantias de visibilidade de
\texttt{synchronized} são tema da seção~\ref{sec:jmm-historia}, e as
alternativas de \texttt{java.util.concurrent}, da
seção~\ref{sec:juc}.
